% SPDX-License-Identifier: CC-BY-SA-4.0
% Session: Form parallelism to concurrency
% Exercise: Asynchrony

\begingroup

\begin{exercice}[Asynchronie]
  \label{exo:introduction/asynchronous}

  On se donne le programme multi-threads suivant : 

  \begin{lstlisting}
    Thread t = new Thread(() -> {
      System.out.print("a");
      new Thread(() -> { System.out.print("b"); }).start();
      System.out.print("c");
    });

    System.out.print("d");
    t.start();
    System.out.print("e");
    t.join();
    System.out.print("f");
  \end{lstlisting}

  \begin{question}
  \item Décrivez la relation ``happens-before'' du programme.
  \item Décrivez toutes les sorties possibles pour le programme, en sachant qu'il est séquentiellement cohérent.
  \item Décrivez la machine à états du programme.
  \end{question}

  \begin{correction}
    Les sorties possibles sont toutes les linéarisations du graphe de causalité, soit les 11 sorties suivantes :
    
    \begin{minipage}{.4\textwidth}
      \begin{tikzpicture}
        \draw (0,0) node {f};
        \draw (1,1) node {c};
        \draw (3,1) node {b};
        \draw (0,2) node {e};
        \draw (2,2) node {a};
        \draw (1,3) node {d};
        
        \draw[-latex] (0.8,2.8) -- (0.2,2.2);
        \draw[-latex] (1.2,2.8) -- (1.8,2.2);
        \draw[-latex] (0,1.8) -- (0,0.2);
        \draw[-latex] (1.8,1.8) -- (1.2,1.2);
        \draw[-latex] (2.2,1.8) -- (2.8,1.2);
        \draw[-latex] (0.8,0.8) -- (0.2,0.2);
      \end{tikzpicture}
    \end{minipage}
    \begin{minipage}{.4\textwidth}
      \begin{itemize}
      \item deacfb
      \item deacbf
      \item deabcf
      \item daecfb
      \item daecbf
      \item daebcf
      \item dabecf
      \item dacefb
      \item dacebf
      \item dacbef
      \item dabcef
      \end{itemize}
    \end{minipage}
  \end{correction}
  
\end{exercice}

\endgroup
\endinput
